\documentclass[11pt,a4paper]{article}
\usepackage{fontspec,amsmath,amssymb,amsthm,microtype}
\setmainfont{lmroman10-regular.otf}[BoldFont=lmroman10-bold.otf,ItalicFont=lmroman10-italic.otf,BoldItalicFont=lmroman10-bolditalic.otf]
\usepackage[margin=25mm]{geometry}
\usepackage[unicode,hidelinks]{hyperref}
\hypersetup{pdfauthor={LangevinLean Project}}
\setlength{\emergencystretch}{3em}
\setlength{\parskip}{0.4em}
\setlength{\parindent}{0pt}
\allowdisplaybreaks[2]
\raggedbottom
\title{\texorpdfstring{Actual-law constant and varying-step averaged LMC bounds}{Actual-law constant and varying-step averaged LMC bounds}}
\author{LangevinLean Project}
\date{}
\begin{document}\maketitle
Draft mathematical exposition; these notes may contain mistakes.


The scalar recurrences are now instantiated by the actual Gaussian algorithm. Its marginal probability laws and finite moments are derived at every step. The exact constant-step exponential bound and the complete variable-step proposition, including its warmup and accuracy ceiling, follow with the paper’s constants.



The actual Gaussian LMC laws satisfy the constant-step bound, including arbitrary initial probability laws in extended transport, when \(0\le h\le1/L\):


\[W_2^2(\nu_k,\pi)\le e^{-2mkh}W_2^2(\nu_0,\pi)+\frac{(M_{\rm av}+m)hp}{2m}.\]


For Proposition 1 put \(\kappa=L/m\), \(a=M_{\rm av}/m+1\), and assume \(W_2^2(\nu_0,\pi)\le\varepsilon_0^2p/m\). With \(\log_+u=\max\{\log u,0\}\), the exact schedule and error envelope are


\[\begin{aligned}k_0&=\left\lceil\frac\kappa2\log_+\!\left(\frac{2\kappa\varepsilon_0^2}{a}\right)\right\rceil,\qquad h_k=\frac1{L+m(k-k_0)_+},\\ W_2^2(\nu_k,\pi)&\le\frac{a}{\kappa+k-k_0}\frac pm\qquad(k\ge k_0).\end{aligned}\]
\section{\texorpdfstring{Define actual marginal laws and propagate finite moments}{Define actual marginal laws and propagate finite moments}}


The recursive law at step \(k+1\) is the independent Gaussian Euler transition at the deterministic step \(h(k)\), applied to the current law. The zero and successor identities expose this recursion. A pushforward of a product of probability laws remains a probability law, proving normalization at every \(k\). For each nonnegative step, the already constructed \(L^{2}\) Brownian realization gives finite output transport whenever initial transport is finite; the zero-step identity covers \(h=0\). Induction gives finite moments for every law of an arbitrary nonnegative deterministic schedule.



Earlier results used here: \cite[Section~3.17.3]{LangevinLean2026ULASC}; \cite[Section~3.26.3]{LangevinLean2026ULASC}; \cite[Section~3.36.7]{LangevinLean2026ULASC}; \cite[Section~3.45.6]{LangevinLean2026ULASC}.

\section{\texorpdfstring{Remove the nonpositive score correction and iterate a constant step}{Remove the nonpositive score correction and iterate a constant step}}


When \(hL\le1\), the score correction coefficient is nonpositive and the score second moment is nonnegative, so the exact one-step theorem implies the two-term affine recurrence. A constant schedule agrees by induction with iteration of the actual Gaussian transition. Strong convexity and coordinate smoothness imply \(m\le M_{\mathrm{av}}\le L\), hence \(0\le mh\le 1\). Apply the scalar recurrence theorem with \(t=mh\) and forcing coefficient \((M_{\mathrm{av}}+m)p/m^2\). Its exponential contraction becomes \(\exp(-2mkh)\), and its fixed point becomes exactly \((M_{\mathrm{av}}+m)hp/(2m)\). This proves Theorem 1’s global claim for actual LMC laws.



Earlier results used here: \cite[Section~3.17.3]{LangevinLean2026ULASC}; \cite[Section~3.17.4]{LangevinLean2026ULASC}; \cite[Section~3.26.3]{LangevinLean2026ULASC}; \cite[Section~3.36.7]{LangevinLean2026ULASC}.

\section{\texorpdfstring{Normalize the variable schedule without changing its constants}{Normalize the variable schedule without changing its constants}}


Multiplying the proved stable one-step inequality by \(m/p\) gives the dimensionless recurrence for \(r(k)=m\) \(W_{2}^{2}/p\) with \(a=M_{\mathrm{av}}/m+1\). The schedule is initially written as the dimensionless schedule divided by \(m\); elementary division shows it is exactly \(1/(L+m(k-k_{0})_{+})\), as displayed in Proposition 1. The scalar warmup proof is extended to zero initial error: then \(k_{0}=0\) and the harmonic tail applies immediately. For nonzero initial error, its absolute value has the same square and the same warmup, so the existing positive-parameter scalar theorem applies without adding a sign restriction.



Earlier results used here: \cite[Section~3.17.3]{LangevinLean2026ULASC}; \cite[Section~3.17.4]{LangevinLean2026ULASC}; \cite[Section~3.26.3]{LangevinLean2026ULASC}; \cite[Section~3.36.7]{LangevinLean2026ULASC}.

\section{\texorpdfstring{Prove the entire variable-step proposition for the actual chain}{Prove the entire variable-step proposition for the actual chain}}


The coordinate comparison gives \(\kappa =L/m\ge 1\) and \(a>0\). The chosen schedule is positive and satisfies \(h(k)L\le 1\) at every step. Use the actual-law normalized estimate and the derived finite moments to discharge every recurrence premise. The warmup and harmonic-tail theorem then gives \(r(k)\le a/(\kappa +k-k_{0})\) for all \(k\ge k_{0}\). Multiply back by \(p/m\) to recover the exact paper envelope. Finally, the explicit positive-part ceiling defining the sufficient iteration count makes this envelope at most \(\varepsilon ^{2}\); multiplying by \(p/m\) proves the stated sufficient iteration threshold for the actual algorithm.



Earlier results used here: \cite[Section~3.17.3]{LangevinLean2026ULASC}; \cite[Section~3.17.4]{LangevinLean2026ULASC}; \cite[Section~3.26.3]{LangevinLean2026ULASC}.

\section{\texorpdfstring{Remove separate initial-moment assumptions from the global statements}{Remove separate initial-moment assumptions from the global statements}}


The exponential contraction factor is strictly positive. If the initial extended Wasserstein distance is infinite, the extended global constant-step inequality has infinite right side and follows immediately. If it is finite, the proved real-valued estimate and finite-output transport convert back to the extended nonnegative inequality. For the variable-step proposition, its stated finite initial squared-distance bound already forces finite initial distance; taking real values gives exactly the initial-error inequality used above. Applying the actual-law variable envelope and accuracy theorems therefore requires no independent finite-moment premise.



Earlier results used here: \cite[Section~3.17.3]{LangevinLean2026ULASC}; \cite[Section~3.17.4]{LangevinLean2026ULASC}; \cite[Section~3.26.3]{LangevinLean2026ULASC}; \cite[Section~3.36.7]{LangevinLean2026ULASC}.

\begin{thebibliography}{ULA-LSI}
\bibitem[ULA-SC]{LangevinLean2026ULASC}
LangevinLean Project. \emph{Unadjusted Langevin under Strong Convexity}. LL-01, 2026. Draft report, content revision 263fb89cdbf0. \url{https://odakyildiz.github.io/LangevinLean/docs/pdf/UnadjustedLangevin/FullProof.pdf}.
\end{thebibliography}

\end{document}
